\scp{Evidence of on-going movement and contact}{02-02}

The information in this section is not intended to be exhaustive,
but merely a sample to illustrate some of the complexities that
need to be understood in the region.
The information is broadly organised by topics,
and within each subsection there is an attempt to be geographical
and chronological where possible.\footnote{
		A reviewer suggested that this section be reorganised chronologically.
		However, there are too many unknowns
		and disagreements among the experts about when many of these
		issues involving contact in the region began,
		to be able to do so coherently.
		The earliest spice trade had a different sociolinguistic
		impact on the region than the much later European quests -- but
		they both involve the spice trade,
		so there is some necessary redundancy in our discussion.
		While there is clear evidence that trade in cloves spans at l\ili{east}
		two millennia, there is some evidence pushing it back to almost four millennia.
		Sandalwood trade may span two millennia or more.
		There could have been some overlap with the Lapita trade.
		How does one sequence those unknowns chronologically?
		It is not known how far back bird-of-paradise trade goes.
		It is not known when the tripang (sea cucumber) trade began,
		and there is also disagreement.
		Technologically it could have been going on for centuries before
		the 300 year-old threshold touted by some.
		The presence of loanwords such as ‘money’
		and ‘chain’ from Austronesian languages into Australian Aboriginal
		languages does not allow us to identufy when the trade began,
		but is merely evidence that the contact continued into
		the historical present.}

\sscp{Trade networks that span centuries and even millennia}{02-02-01}
There has been much trade through the region,
it has involved people from outside the region,
and it has been happening for a long time -- some of it spanning two millennia or more.

The Lapita cultural complex associated with Oceanic languages in the Pacific involved
trade spanning thousands of kilometres
and several centuries, perhaps beginning around 3470--3250 cal.
BP \citep[95]{SpechtTerrell2014}.
The Lapita trade involving pottery
and obsidian in near Oceania also had links back to Borneo
and to the northern Philippines -- and was thought to have links to Timor and Maluku.
\citet{SpechtTerrell2014} explain these Lapita connections as follows:

\begin{quote}
	``Interisland movement of goods also has a long history within ISEA,
	as obsidian in late Pleistocene
	and mid-Holocene contexts on Timor,
	Borneo, and the \ili{Talaud} Islands almost certainly derived from off-island sources
	\citep{AmbroseReepmeyer2009,Reepmeyer2011,Reepmeyer2011}.
	Contacts with the New Guinea region in the early
	and mid-Holocene are attested by the recovery of marsupial faunal remains in prepottery
	contexts in Timor and the \ili{Talaud} and northern Molucca Islands
	\citep{FlanneryNitihaminoto1995,FlanneryBalasubramaniam1998}.

	The first clear evidence for contact between ISEA
	and the [Bismarck] Archipelago comes from the transfer of pottery technology,
	followed by the commensal animals discussed above.
	In the reverse direction,
	New Britain obsidian occurs at Bukit Tengkorak on Borneo in contexts contemporary
	with Middle Lapita
	\citep{BellwoodKoon1989,Bird1989}.''
	\citep[116--117]{SpechtTerrell2014}

	``The Borneo connection with the [Bismarck] Archipelago continued after Lapita.
	Two undated double-spouted pots found in the \ili{Manus} Islands show ‘‘a singular similarity’’
	\citep[p. 28]{Kennedy1982} to vessels of uncertain age in Niah caves in Sarawak \citep{Harrison1971},
	and a piece of bronze of probable ISEA origin found on
	\ili{Manus} is assigned a post-Lapita date \citep{Ambrose1988}.
	In summary, there is evidence in pre-Lapita times of long-term interaction between
	ISEA and New Guinea and between New Guinea
	and the Archipelago, but not for direct interaction between ISEA and the Archipelago.''
	\citep[117]{SpechtTerrell2014} 
\end{quote}


It is important to understand that these are pointing to \it{trade} links involving
travel in two directions spanning a long period of time.
This provides early direct evidence of things like obsidian associated with the Lapita
cultural complex, sourced in the Bismarck Archipelago north
and \ili{east} of New Guinea,
having been brought westward to Borneo,
going against the flow of the singular idea of a monolithic
eastward migration by Austronesian-speaking peoples.
And this trade involved Sulawesi and Borneo.
There appears to have been travel back-and-forth,
and contact maintained over a significant period of time.
The trade may not have involved end-to-end trade routes,
but rather worked through intermediate links -- much like the spread
of {\#}muku ‘banana’ (see \srf{02-03-01-07}).

Earlier ideas that these Lapita links somehow involved north Maluku are being questioned by some.
\citet[219]{Bellwood2019a} states:

\begin{quote}
	``Actual cloves have been found in early to middle first millennium AD contexts at
	Segaran IIA in West Java
	and Mantai in Sri Lanka \citep{ManguinIndradjaja2011}.
	Furthermore, Indian pottery in quite large quantities is now well known from the
	sites of Sembiran and Pacung in northern Bali,
	dating to circa 2000 years ago or just before
	\citep{ArdikaBellwood1991,ArdikaYuliathi1997,Calo2015}.
	Neither ancient cloves nor Indian pottery have been excavated in the Moluccas,
	but it seems relatively certain that Indian contact could \it{potentially} have reached
	the Moluccas by 2000 years ago.
	So could, and did,
	the technological developments that are widely associated by archaeologists,
	linguists, and historians with a combination of Indian,
	Mainland South\ili{east} Asian, early Malayo-Chamic,
	and Chinese trading acumen.
	These include metal items
	and metallurgy, Dong Son drums, glass beads,
	nephrite ornaments (although these are not yet found in the Moluccas),
	and probably also some new Y-chromosome haplogroups \citep{Kusuma2016}.''
	\citep[219]{Bellwood2019a} [emphasis in the original]
\end{quote}

The later so-called “spice orgy” (e.g. \sqcitet{Wright2007}{Wright2007} among others)
associated with Europe
and eastern Indonesia has impacted the history of the world,
including Columbus bumping into America on his way to find a back door to the Spice Islands of Maluku.
It also has a role in the histories of the Netherlands,
Ceylon, Capetown, Manhattan, Suriname,
and economies of islands such as Granada in the Carribean,
and Zanzibar in Africa.
(See \sqcitet{Milton1999}{Milton1999} as one of many sources.)

Cloves (\it{Syzygium aromaticum} [formerly \it{Eugenia caryophyllatta}]),
nutmeg and mace (\it{Myristica fragrans}) are native to northern
and central Maluku respectively (see \frf{02-07} to \frf{02-10}).
Sandalwood (\it{Santalum album Linn}.) is thought to be indigenous to Timor
and Sumba and the Indian subcontinent.
Trade in these items goes back hundreds,
and even thousands of years,
long before European presence in the region.

\begin{figure}
	\centering
	\frame{\includegraphics[width=0.3\columnwidth]{02-07.jpg}}
	\fcp{Drying cloves for shipping}{02-07}
\end{figure}

Although the exact trade routes are not known,
\citet[67]{Lapian1965} reports finds of cloves in present day Syria in the remains of a
pantry room in a common household,
dating to around 1,700 BCE.\footnote{
		While the botanical identity of the Terqa cloves
		has been questioned, the issue has not been settled,
		nor does it detract from other evidence of clove trade spanning over 2,000 years
		(nearly 4,000 years if the Terqa cloves are indeed cloves).}
\citet[1--2]{Andaya1993} notes
that some of the earliest written references to cloves include a Chinese source during
the Han dynasty (206 BC-222 AD).
Pliny the Elder describes cloves reaching the Roman world by the first century AD.
\citet[67]{Lapian1965} finds clear evidence that cloves had reached Europe by the 4th century,
when Silvester, the Bishop of Rome from 314--335,
received a gift of 150 pounds of cloves.
\citet[78]{Lapian1965} also notes that before Europeans gained direct access to the source
of the spices in Maluku in 1512,
that Asian traders brought the spices from Maluku to the western part of the archipelago
where they were then purchased by Arab,
Chinese and other foreign traders
and eventually made their way to China,
the Middle East, and Europe via Venice.
There is some evidence that many of these “Asian” traders were from Java
and \ili{Makasar} and were associated with the coming of trade \ili{Malay} into the eastern regions,
already functioning as a language of wider communication by the time of the arrival
of the Europeans \citep{Grimes1991,Grimes1994}.

%%%%%%%%%%%%%%%%%%%%%%%%%%%%%%%%%%%%%%%%%%%%%%%%%%%%%%%%%%%%%%%%%%%%%%%&%
%%                                                                     %%
%% NOTE FOR TYPESETTING: I have added some manual spacing to have      %%
%% the pictures not too close together, but also not too far apart     %%
%% and to have them appear on the same page as the text referencing    %%
%% them. The sizes have also been reduced (they were too big in        %%
%% the proofreading version.                                           %%
%%                                                                     %%
%% Please consult before editing how these pictures have been typeset! %%
%% Changing the sizes could have bad flow on effects for pagination.   %%
%%                                                                     %%
%%%%%%%%%%%%%%%%%%%%%%%%%%%%%%%%%%%%%%%%%%%%%%%%%%%%%%%%%%%%%%%%%%%%%%%%%

\begin{figure}\centering
  \subfigure[Fresh picked cloves \it{(Syzygium aromaticum)}]{\label{fig:02-08}
    \frame{\includegraphics[width=0.31\textwidth]{02-08.jpg}}}\hspace{2mm}
  \subfigure[Preparing nutmeg and mace for drying]{\label{fig:02-09}
    \frame{\includegraphics[width=0.31\textwidth]{02-09.jpg}}}\hspace{2mm}
  \subfigure[Drying nutmeg \it{(Myristica fragrans)}]{\label{fig:02-10}
    \frame{\includegraphics[width=0.31\textwidth]{02-10.jpg}}}
\fcp{Spices native to eastern Indonesia}{02-08--02-10}
\end{figure}

%\begin{figure}
  %\subfigure[Bronze Dong-Song drum from Vietnam used as bridewealth in \ili{Alor}]{\label{fig:02-11}
		%\includegraphics[width=0.25\textwidth]{02-11.jpg}}\hspace{10mm}
  %\subfigure[One of many  varieties of tripang \it{(Holothuroidea sp.)}]{\label{fig:02-12}
		%\includegraphics[width=0.25\textwidth]{02-12.jpg}}\hspace{10mm}
	%\subfigure[Sign at ``Macassan   Beach'' in north\ili{east}   Arnhem Land, northern Australia]{ \label{fig:02-13}
		%\includegraphics[width=0.25\textwidth]{02-13.jpg}}
%\caption{Long-distance trade\label{fig:02-11--02-13}}
%\end{figure}


%\begin{figure}
    %\subfigure[Fresh picked cloves  \it{(Syzygium aromaticum)}]{\label{fig:02-08}
		%\frame{\includegraphics[width=0.25\textwidth]{02-08.jpg}}
%}
    %\subfigure[Preparing nutmeg  and mace for drying]{\label{fig:02-09}
		%\frame{\includegraphics[width=0.25\textwidth]{02-09.jpg}}
    %}
    %\subfigure[Drying nutmeg   \it{(Myristica fragrans)}]{\label{fig:02-10}
		%\frame{\includegraphics[width=0.25\columnwidth]{02-10.jpg}}
%}
	%\fcp{Spices native to eastern Indonesia}{02-08--02-10}
%\end{figure}

Timor and Sumba figured in the sandalwood trade.
\citet[1]{Hagerdal2012} gives the first historical
mention of the sandalwood trade on Timor
from around 1350 \tsc{ce} in Chinese sources.
Since contact with the Indian subcontinent dates back to the
first and second centuries \tsc{ce},
it is possible that sandalwood trade in Timor
and Sumba could pre-date the first mention in Chinese sources
by several centuries.

There were early trade networks linking Vietnam
and Thailand to the islands of eastern Indonesia.
\citet[17]{Klamer2019} says, 

\begin{quote}
``During the Metal Age ({\tl}3,000--1,500 BP),
traces of trade networks become more visible.
For example, Galipaud (in press) mentions three types of links: one linking south
Vietnam and Thailand and the Indonesian islands,
the other linking northwest Vietnam,
Laos and Yunnan to Sumatra,
eastern Indonesia and New Guinea, and the third linking south Vietnam with the Philippines and Serawak.
South Vietnam and Thailand are linked with the Indonesian islands way into the \ili{east},
to the Lesser Sunda islands Flores,
Lembata, \ili{Pantar}, and \ili{Alor} by the existence of Dong-Song drums
and the existence of cemeteries with burials where bodies were put in certain types of jars.
The transport of some Dong-Song drums into ISEA is one of the clearest traces of
a regional connection between the SE Asia mainland
and the islands and is dated after the dispersal of MP [Malayo-Polynesian] languages.
The cemeteries with urn containers are always coastal
and therefore suggestive of mobile communities connected by a vast network of social
relations extending possibly as far as the South\ili{east} Asian mainland.
In a recent excavation of the cemetery of Pain Haka in eastern Flores
\citep{GalipaudBuckley2016}, the skeletons were dated
and revealed that urn burials were already in use by 2,800 BP in Pain Haka,
and that the cemetery remained in use for more than 500 years,
without any evidence of imported metal.
Pain Haka excavations provide clear evidence of burial practices similar to those
documented in other parts of SE Asia.
Chief among these is the use of pottery jars alongside other forms of container for
the interment of the dead.
The dating of the site combined with the fact that this burial practice is present
over such a wide geographic area suggests a widespread belief system during the Neolithic
period across much of SE Asia.
In other words, burial traditions linking ISEA with the South\ili{east} Asian mainland
started well before the Metal Age,
i.e., before the dispersal of MP languages.'' \citep[17]{Klamer2019}
\end{quote}

Tripang trade (sea slug,
sea cucumber; \it{Holothuroidea spp}.; see \frf{02-12})
and \it{Trochus} (top shells) used in making buttons
and jewellery also brought traders from many other parts of Asia into our target region.
They were also a catalyst for early contact with northern Australia (particularly
Arnhem Land) by seafarers from what is now central (including Makassar)
and southeastern Indonesia (including the Lesser Sundas and Maluku).
Scores of boats made the yearly trek to northern Australia,
planted tamarind trees, and left archaeological sites,
pottery, words, and part of their gene pool.
The term “Macassan” in northern Australia is a grab bag to refer to anyone from the
Indonesian archipelago, and is not limited to Makassar.
Hundreds of loanwords from Austronesian languages have been identified in Aboriginal
languages of northern Australia \citep{WalkerZorc1981,Evans1992}.
Grimes has been told stories by Aboriginal people in both \ili{east}
and west Arnhem Land of their grandfathers
and great-grandfathers accompanying boats on the return journeys northward
and marrying local women in central
and eastern Indonesia, some returning to their homelands after decades away.

Yolŋu Matha words like orthographic ‹\it{dhimurru}› [ˈd̪imuru] ‘\ili{east} wind,
\ili{east} monsoon’ and ‹kapala› [ˈkapala] ‘ship’ with initial stress reflect the
telltale extrametrical echo vowels on roots ending with \it{r,
l, s} that narrow the probable source languages to Makassar,
Konjo, or \ili{Selayar} in South Sulawesi \citep[96]{GrimesGrimes1987}.
Words like ‹djalataŋ› [ˈʤalataŋ] ‘south’
and ‹rranhdhiŋ› [ˈran̪d̪iŋ] ‘chain’ are probably from some form of \ili{Malay} (cf.
\ili{Malay} \it{səlatan} ‘south’; \it{rantai} ‘chain’).
But ‹bärra› [ˈbara] {\tl} [ˈbaːra] ‘west wind,
west monsoon’ and ‹djarra› [ˈʤara] ‘horse’ could be from many possible sources;
for example many languages in the Lesser Sundas have \it{jara} [ˈʤara] ‘horse’,
whereas the final nasal of Makassar \it{ʤarraŋ} ‘horse’ is commonly preserved in
Yolŋu Matha languages and so is unlikely to be the source.

%%%%%%%%%%%%%%%%%%%%%%%%%%%%%%%%%%%%%%%%%%%%%%%%%%%%%%%%%%%%%%%%%%%%%%%%
%%
%% NOTE FOR TYPESETTING: I have added some manual spacing to have
%% the pictures not too close together, but also not too far apart
%% and to have them appear on the same page as the text referencing
%% them. The sizes have also been reduced (they were too big in
%% the proofreading version.
%%
%% Please consult before editing how these pictures have been typeset!
%% Changing things could have bad flow on effects for pagination.
%%
%%%%%%%%%%%%%%%%%%%%%%%%%%%%%%%%%%%%%%%%%%%%%%%%%%%%%%%%%%%%%%%%%%%%%%%%

\begin{figure}\centering
  \subfigure[Bronze Dong-Song drum from Vietnam used as bridewealth in \ili{Alor}]{\label{fig:02-11}
		\includegraphics[width=0.25\textwidth]{02-11.jpg}}\hspace{10mm}
  \subfigure[One of many  varieties of tripang \it{(Holothuroidea sp.)}]{\label{fig:02-12}
		\includegraphics[width=0.25\textwidth]{02-12.jpg}}\hspace{10mm}
	\subfigure[Sign at ``Macassan   Beach'' in north\ili{east} Arnhem Land, Australia]{ \label{fig:02-13}
		\frame{\includegraphics[width=0.25\textwidth]{02-13.jpg}}}
\caption{Long-distance trade\label{fig:02-11--02-13}}
\end{figure}

%\begin{figure}
    %\subfigure[Bronze Dong-Song drum from Vietnam used as bridewealth in \ili{Alor}]{\label{fig:02-11}
	    %\frame{\includegraphics[width=0.3\textwidth]{02-11.jpg}}
    %}
  %\subfigure[One of many  varieties of tripang \it{(Holothuroidea sp.)}]{\label{fig:02-12}
	%\frame{\includegraphics[width=0.3\textwidth]{02-12.jpg}}
	%}
    %\subfigure[Sign at ``Macassan   Beach'' in north\ili{east}   Arnhem Land, northern Australia]{ \label{fig:02-13}
	%\frame{\includegraphics[width=0.3\columnwidth]{02-13.jpg}}
	%}
%\caption{Long-distance trade\label{fig:02-11--02-13}}
%\end{figure}

In a recent article,
\citet{Schapper2022} has identified a handful loanwords in indigenous languages of
northern Australia which originate from languages within our target region: terms
related to ‘hawksbill turtle’ from reflexes of Proto-\ili{Aru} *kaloba,
terms meaning ‘fishing line’ from \ili{Aru} languages,
‘sea cucumber’ from words related to Proto-Rote-\ili{Meto} *ɓanafi,
and ‘tobacco’ from words related to Proto-Rote-\ili{Meto} *maɗo \citep[9--12]{Schapper2022}.
Schapper posits that such forms indicate that peoples from Wallacea may have undertaken
their own voyages to northern Australia.

%\largerpage
The long-standing seasonal interaction between the Austronesian-speaking peoples
and the people of Arnhem Land is still remembered
and highly valued, even though the Australian government has attempted to cut off
the relationship since Federation in 1901.
For example, on 14 June 2020 the Aboriginal band East Journey performed a concert
singing their songs in the \ili{Gumatj} language at the Bawaka homeland,
%\footnote{
		%See \url{https://www.facebook.com/VisitEastArnhemLand/videos/3812185228851950/?_rdc=1&_rdr}
		%for a recording of the concert.}
as part of a series of Covid-19 online concerts
from several bands in East Arnhemland.
Around 6:58 minutes into the concert they were singing a traditional song about
\it{liya-wayma kamboŋ malogo, liya-wayma kamboŋ dhua},
which in \ili{Malay} would be something like
\it{layar-layar kampuŋ Maluku, layar-layar kampuŋ tua}
‘the sails of Maluku village, the sails of the old village’.
In a contemporary song at around 7:54 minutes into the concert they sing
\it{ga wiripunydja kamboŋ dhua} “and another way/name is \it{kampuŋ tua} [old village]”,
with the next line clarifying “these speak of Bawaka”
(the remote homeland where the concert was filmed).
Several involved with the concert series have names
like Djakapurra (Singapore) and Djawa (Java).
There are songs about places like \it{Gupaŋa} (\ili{Kupang}).\footnote{
		Our appreciation
		to Concert Producer and Sound Engineer Andrew Grimes for narrowing down the phrases for us.
		He gave additional examples,
		but what is presented here suffices to make the point.}
On 9 July 2020 a Facebook
post from Numbulwar, further south in the gulf of Carpentaria,
was seeking assistance for translating a number of names associated with flags
and dances linked to the “Macassans”.
The memory and the significance of the contact is not lost.

While hundreds of loanwords have been identified going from Austronesian languages
into various Australian Aboriginal languages,
nobody seems to have seriously investigated loans going in the other direction.
But until we have better documentation of those languages,
the task would remain difficult.

Coastal west Papua and the \ili{Aru} islands supported a
long-standing trade in bird-of-paradise
specimens and feathers, and other exotic fauna.
This also brought a steady stream of traders from the western parts of the Indonesian
archipelago and further afield \citep{Andaya2017}.

In addition to the well described spice trade routes to \ili{Ternate},
\ili{Ambon} and \ili{Banda}, there was also a trade network from
\ili{Banda} to \ili{Onin} on the west coast of Papua.
The language of \ili{Geser}-\ili{Gorom} has also been used as a regional lingua franca around
eastern Seram and the west coast of Papua for generations
\citep{LoskiLoski1989,WalkerWalker1991,Zobel2014}.
These regional connections are also reflected in clan names on \ili{Ambon},
such as Latupapua ‘king of Papua’ (pers. comm. Zeth Latupapua).
These factors must be kept in mind when unravelling the linguistic
affiliations of \tsc{Seram-Tanimbar-Bomberai} languages.
There are telltale differences in the sound correspondences.
(See \crf{ch:STB}.) 

\citet{Rutter1930} describes seasonal slave raiders travelling from the southern
Philippines through northern Maluku to the Bird’s Head area of Papua,
before continuing on to Timor
and Rote through central Maluku,
and eventually returning north to the Philippines. Traces are widespread.
\citet[74]{Grimes2000a} describes stories on \ili{Buru} about
the \it{geba lano} ‘\it{lano} people’
slave raiders, and the local \ili{Buru} \it{epkitan}
‘fighting champions’ who were to watch for them from
high spots along the coast
and fight to protect their own people when the slave raiders came.
The \it{geba lano} are variously thought in different parts of \ili{Buru} to be from Papua,
from Seram, or from \ili{Tobelo}
and \ili{Galela} (Halmahera, north Maluku; West Papuan languages).
The term \it{lano} may be associated with the \ili{Ternate} (West Papuan language) term
\it{kolano} which refers to a political position
in societies of north Maluku \citep{Andaya1990},
or it may refer to Seram which is often referred to in central Maluku languages as
\it{selano} (as in the clan names Pattiselano or Huliselan, etc; [pers. comm.
Jules Pattiselano; pers. comm. Mus Huliselan]).
But most likely it refers to the seasonal Ilanun
slave raiders from the southern Philippines \citep{Rutter1930}.
Their impact is still felt
and feared after generations.

\begin{figure}
	\centering
	\frame{\includegraphics[width=0.5\columnwidth]{02-14.jpg}}
	\fcp{Traders from south\ili{east} Sulawesi entering \ili{Ambon} harbour}{02-14}
\end{figure}

While the role of the ships
and sailors from \ili{Makasar}
and Java is better known in the historical literature,
inter-island trade in eastern Indonesia into the present has,
for generations, been carried out by sea-going traders variously
associated with the islands of Buton,
Binongko, \ili{Wanci}, \ili{Cia-Cia}, \ili{Tukang Besi}
and others off south\ili{east} Sulawesi (see \frf{02-14}).
There are long-standing settlements from south\ili{east} Sulawesi scattered
throughout the region, including northern coastal Flores,
\ili{Pantar}, Obi, \ili{Buru} and Seram.
For example, \citet{Grimes2000a} notes that many communities along
the western and southern coasts of \ili{Buru} are long-standing settlements totalling
around 12,600 people from south\ili{east} Sulawesi.
Much of the knowledge
and terminology associated with the sea in the \ili{Buru} language,
such as ways to treat a stonefish puncture,
is said to have been learned from these settlers.
There is some evidence from the lexicon
and phonology (such as initial
and medial prenasalised stops) that indicates the recently extinct
language of \ili{Hukumina} in northwest \ili{Buru} was possibly a satellite
community of speakers from Sulawesi \citep{Grimes2010b}.

All of these trade links serve to illustrate the kinds of “contaminants”
that can muddy the waters in trying to unravel the phonological
prehistory of the Austronesian languages in the region.

\sscp{Empire outliers}{02-02-02}
Indian influence via the Majapahit Hindu-Buddhist empire (around
1293--1517) in \ili{east} Java left emissaries in \ili{Ternate}, \ili{Ambon}
(e.g. \it{Rehatta} clan; pers. comm.
Sonja Rehatta) and Timor (e.g.
the \it{Kase} clans; pers. comm.
Simon Kase); a caste system along the southern islands as far
\ili{east} as \ili{Kei}; and kin group names in places such as \ili{Kei}
and \ili{Buru} independently associated with the Majapahit empire
(e.g. \it{Masbait} clan, pers. comm.
Raja Anton Lesnussa). Javanese words are still used as the names
of islands in the region (e.g.
\it{Wetan} ‘\ili{east}’ and \it{Alor} ‘north’), and in names for introduced
fauna (e.g. \ili{Buru} \it{menʤaŋan} ‘deer’).

While there was some trade associated with the Majapahit presence
in eastern Indonesia, it is the empire,
its infrastructure and dominance over vassal states that has
left a lasting legacy in the region.
While Majapahit was the last of the Indic
kingdoms in Indonesia, Indian influence across
Indonesia traces back as early as the
first to second centuries CE.
See also \citet{Hoogervorst2016,Hoogervorst2017,Hoogervorst2023,Hoogervorst2024}
for the influence of South Asian languages in the region.

The people of Sabu often express their perception that their
identity and culture are linked to India.
\citet[199--244]{Fox1977} explores many parallels in the livelihoods
and economies of Rote
and Sabu with southern India
and Sri Lanka, all based around the lontar palm (\it{Borassus flabellifer}).

The political power-house of Gowa (Makassar,
in south Sulawesi, west of the Blust line
and not a “CEMP” language) dominated the Flores
and \ili{Sumbawa} region up until 1727,
leaving linguistic traces of varying degrees \citep[585]{VerheijenGrimes1995}.

After 1727, \ili{Bima} (a “CEMP” language of eastern \ili{Sumbawa}) dominated
the nearby islands around Flores.
This contact may account for some of the similarities behind
the undefined notion of a “\ili{Bima}-Sumba” subgroup,
for which \citet{Blust2008} has not been able to find support
in the historical phonologies.

\citet{vanFraassen1987} and
\citet[15--31, 146, 159--160]{Grimes1993}
describe the history of contact
between \ili{Ternate} and central Maluku (Sula,
\ili{Buru}, \ili{Manipa}, \ili{Kelang}, Boano,
Seram), and the influence from Ternatan ‘governors’ on the political
structures, social titles and language on \ili{Buru}.
The names of several individual Ternatan emissaries are known.
\ili{Ternate} is a Papuan language from north Halmahera (see \frf{02-06}).
\ili{Buru} is an Austronesian language in central Maluku.
The “memory” of the period of Ternatan control on \ili{Buru},
stopped by the Dutch in the mid-1600s,
is still talked about on \ili{Buru} today
and figures prominently in oral histories of the island.
Several of the Ternatan legacy positions/titles are still used
by some kin groups in the mountains.

While a number of \ili{Ternate} loans (a Papuan language) have been
identified in the \ili{Buru} language,
such as \it{matgugul} ‘title of \ili{Ternate} governor’,
\it{soa} ‘clan’ and \it{jou} ‘lord’, the evidence is much more
striking in the Sula Islands north of \ili{Buru}.
Fortgens was the investigator for the Holle lists for \ili{Soboyo},
\ili{Taliabu}, and \ili{Kadai} \citep{Stokhof1980b}.
In the \ili{Soboyo} data Fortgens identified 52 loans from \ili{Ternate}
and 4 from Java.
In the \ili{Kadai} data he identified 43 loans from \ili{Ternate}
and 3 from Java.
\citet{Collins1982c} has multiple references of \ili{Ternate} influence
and presence throughout Seram
and the islands in the \ili{Manipa} Straits,
as well as noting the presence of several \ili{Tobelo}
and \ili{Galela} speaking villages (both Papuan languages from north
Halmahera) around coastal western and northern Seram.
\ili{Ambon} \ili{Malay} \it{kusu} ‘cuscus’ is very likely a loan from a \tsc{North
Halmahera} language, such as \ili{Ternate} \it{kuso} ‘cuscus’,
particularly since no other language in central Maluku has a
form remotely similar in form (see \crf{ch:Mar}, particularly \srf{14-03-02}).

\sscp{Colonial greed and forced resettlements}{02-02-03}
The events described below
and their impact on language have been written about in more detail by
\citet{Grimes1994,Collins2003b,MusgraveEwing2006,Grimes2010b}.

Nutmeg (\it{Myristica fragrans}) is native to the \ili{Banda} Islands.
In 1621, Dutch \it{Vereenigde Oost-Indische Compagnie} (VOC)
forces decimated perhaps 90{\%} of the native population in order
to gain control of the islands
and the trade in nutmeg.
Remnants of the \ili{Banda}-speaking population fled to the south\ili{east}
to the \ili{Kei} islands,
and set up the refugee communities of \ili{Banda}-Elat
and \ili{Banda}-Eli that remain today,
albeit with a very different pattern of language contact than they had formerly.

The Hoamoal wars of 1651--1656 in central Maluku were an effort
by the Dutch VOC to gain a monopoly in the spice trade,
and limit planting to keep prices high.
The VOC took leaders hostage from many coastal communities,
and forcibly resettled whole or partial language communities
around their forts, creating multiethnic, multilingual communities around several
points of concentration (e.g.
\ili{Batumerah}, \ili{Manipa}, \ili{Kayeli}; see \frf{02-15}).
\citet[17]{vanFraassen1983} describes it as follows:

\begin{quote}
		``The war ignited by the 1651 rebellion [against the Dutch] continued until 1656.
		The rebels received Macassarese support,
		and the war was in no way restricted to Hoamoal [west Seram].
		Arnold de Vlamingh van Oudtshoorn,
		the Dutch commander-in-chief in this war,
		was also launching attacks on the VOC’s opponents elsewhere,
		among other places in \ili{Buru} and \ili{east} Seram.
		[{\ldots}] The Ambonese region was formally removed from Ternatan
		control and the institution of the Ternatan governorship was abolished.
		Hoamoal, [{\ldots}] was completely depopulated [{\ldots}] The
		chiefs of Hoamoal had a place of residence assigned to them in
		\ili{Batumerah}, in the vicinity of the VOC’s chief fortress in \ili{Ambon}.
		The population of the islands of Boano,
		\ili{Kelang} and \ili{Ambelau} was transferred to the island of \ili{Manipa},
		where the Company had a small fort.
		All the Muslim chiefs of the coastal areas of \ili{Buru} were obliged
		to settle in the neighbourhood of the Company’s fort at \ili{Kayeli}.
		[{\ldots}] The evacuated areas were systematically destroyed
		and rendered unfit for reoccupation.
		Contracts were concluded with the chiefs of \ili{Buru}
		and North and East Seram in which the latter conceded themselves
		to be subordinate to the VOC; promised to entertain no relations
		with other nations or rulers,
		to keep out all Macassarese,
		Malays, and Javanese.'' \citep[17]{vanFraassen1983}
\end{quote}

\begin{figure}
	\centering
	\frame{\includegraphics[width=0.45\columnwidth]{02-15.jpg}}
	\fcp{Entrance to VOC fort at Kayeli}{02-15}
\end{figure}

Note the repeated references to alliances with “Macassarese, Malays, and Javanese”.
The fallout from these events continues to shape the sociolinguistic
landscape of central Maluku today,
nearly 400 years later.
The last speakers of several of these transplanted language communities
have just died out in the last three or four decades,
and many of their descendants have completed the shift to other
languages such as \ili{Ambon} \ili{Malay}.

\sscp{Internal migrations}{02-02-04}
Oral histories throughout the region indicate continuous
cross-fertilisation with regional migrations both pre- and post-European contact.
For example, local origin myths tie \ili{Ambelau} (south of \ili{Buru}) together
with \ili{Nusalaut} (south of Seram).
There are clans in several language groups around Timor that
claim their origins trace to Seram,\footnote{
Seram written with
a final m{\#} is from a Portuguese orthographic tradition that
writes final nasals as ‹m›.
Locally in central Maluku it is pronounced [ˈseran] or more frequently [ˈseraŋ].
\ili{Gorom} is similarly pronounced [ˈgoroŋ] around central Maluku.
On Timor, orthographic ‹Tetum› is pronounced [ˈtetʊn].
Compare the pronunciation of Portuguese orthographic \it{Belém}
[beˈleŋ], \it{sim} ‘yes’, etc., which are also not bilabial.}
such as the \it{Seran} clans in the \ili{Tetun}-speaking region (pers. comm. E. Seran;
pers. comm. Father Dr. M.
Bria Seran), and another in the \ili{Helong}-speaking region.
It remains to be seen if these are merely folk etymologies or
if there is more to them.
Given the complex history of the region,
it is not beyond possibility.

In most regions of eastern Indonesia there are origin groups
(kin groups, clans) whose origin stories are local.
They believe they are autochthonous,
originating from the island where they now are.
The origin stories of other groups have a second type of people
originating elsewhere and migrating in the distant past to their current locations.
Yet these latter second-wave groups are also fully assimilated
into the local social structures
and languages, and may even be given the role of a leader or
king \citep{Fox1995,Grimes2006}.

For example, \citet[77--84]{Grimes1993} describes a number of \it{autochthonous}
origin groups on \ili{Buru} whose stories have their founding ancestor
or ancestress appearing out of a rock,
a tree or a stream on the island.
The closer to the headwaters,
the more senior in precedence in relation to other origin groups
(but since each group governs itself,
precedence does not give them any governing rights over another
group; it only provides a script for how elder-younger social relationships should interact).
The further downstream the origin story,
the more junior in precedence.
The place where the founding ancestor built the first house is
known as the \it{tean elen} ‘place of planting the house pole into the ground’.

Other \it{non-autochthonous} origin groups came from elsewhere,
settled on \ili{Buru}, learned the language,
intermarried, and became assimilated into \ili{Buru} culture.
The location where their founding ancestor or ancestress first
set foot on \ili{Buru} is known from their origin stories,
and is called the \it{waga eʔnohon elen} ‘place of disembarkation
from the boat’ (from /en-toho-n/ with regular participial metathesis
from \it{toho} ‘descend, disembark’ < PMP *tuRun ‘descend’),
or an \it{esnaba} ‘disembarkation’.
For example, the \it{Mual} kin group attributes its origins to
the arrival of their ancestress,
Bokisraja, from the Hoamoal peninsula of Seram (note the phonological
similarity between \it{Mual} and the final part of Hoa\emph{moal}).
Their \it{esnaba} is near the village of Masnana.
The Nalbessy kin group \it{esnaba} is near the village of Wae Hotong.
Neither of those groups live near their \it{esnaba}.
Other kin groups claim to have come from other unnamed islands.
But power is tied to place,
so in 1990 when a non-autochthonous kin group that had migrated
to lake \it{Rana} (< PMP *danaw ‘lake’) in the centre of the
island had a ‘bad death (many people die in a short span of time)’,
they abandoned their village at \ili{Rana} where they had lived for
many generations and migrated four days walk to the northwest coast of \ili{Buru} to
reconnect with the power associated with their ancestral \it{waga
eʔnohon elen} near Bara to try
and restore harmony and prosperity.

\citet{Grimes1993} also discusses how these \it{non-autochthonous}
kin groups are categorised as \it{geb fuk Buru} ‘people of \ili{Buru}
island’, and not as \it{geba man lau tau-n} ‘outsiders’ (lit: people whose
origins flow from the sea),
such as the migrants from Sula
and Sulawesi that have been on \ili{Buru} for 12--15+ generations.
Similar patterns are found throughout the region.

The entire north coast of \ili{Buru} has been settled for 15+ generations
by around 12,000--13,000 people from the Sula islands to the
north, completely dominating speakers
of the \ili{Lisela} language on \ili{Buru} \citep{Grimes2000a}.
This has been largely ignored in attempts to establish a linguistic
link between \tsc{Buru} and \tsc{Sula}, which share only about 27{\%}
of their vocabulary
(even less between \tsc{Buru}
and \tsc{Taliabu}), and have very different historical phonologies,
although they do share PMP *R > \it{h}.
(See chapters \ref{ch:Tal} and \ref{ch:SB}.) 
The west coast of \ili{Buru} has been settled by migrants from south\ili{east}
Sulawesi since recorded history.
Their long-standing dominance has resulted in loans into the
\ili{Buru} language, and in things like the corruption of place names such as \ili{Buru}
\it{Kak Pala Madat} (the name of the highest mountain in central
and southern Maluku at 2,428m; \it{kak < kaku} ‘mountain,
ridge’ also found in other place names such as \it{Kak Tuan,
Kak Hanga}, etc.) to \ili{Sanskrit}-\ili{Malay} \it{Kapala Madan},
which recently became the formal government name for the new
regency in western \ili{Buru} where there are very few native \ili{Buru}
people to stand up for their heritage or tell their histories.
Several village names on coastal western \ili{Buru} have been similarly
corrupted or have outside names.
This has not happened nearly so much in areas where these immigrants
are not so dominant.

Volcanic activity around 1976 near the southern Maluku islands
of \ili{Teun}, \ili{Nila}, and \ili{Serua} (\tsc{Timor-Babar} languages)
triggered a government-sponsored
forced resettlement of the speakers of those three languages
to central Seram, hundreds of kilometres to the north,
which has severely impacted the vitality of these languages \citep{Edwards2025}.

What is described above is only a sample for illustrative purposes,
not an exhaustive list.
Some historical linguistic studies in eastern Indonesia have
overlooked these events that have triggered contact-induced language
change in the region -- including from other Austronesian languages.
Such contact issues must be sorted out before comparative work
based on word lists in isolation can be substantiated.


